\documentclass[12pt]{article}
\usepackage[utf8]{inputenc}
\usepackage[spanish]{babel}
\usepackage{amsmath}
\usepackage{amsfonts}
\usepackage{amssymb}
\usepackage{geometry}
\geometry{letterpaper, margin=2.5cm}

\begin{document}

\begin{center}
    \Large \textbf{Evaluación de Examen 1: Unidad I} \\
    \large Física para Ciencias de la Salud (005-1713) \\
    \vspace{0.5cm}
    \textbf{Estudiante:} Udrannys Rojas \quad \textbf{C.I.:} 32.824.239
    \vspace{0.3cm} \\
    \large \textbf{Calificación Final: 6/10 puntos}
\end{center}

\vspace{0.5cm}

\section*{Análisis de Resultados}

\subsection*{1. Ángulo plano (Conversión a radianes)}
\textbf{Procedimiento y resultado:} Plantea correctamente una regla de tres con la equivalencia $180^\circ \rightarrow \pi$ rad. Desarrolla la multiplicación y simplifica de forma directa a la fracción $\frac{5}{12}\pi$ rad. Procedimiento claro y correcto. \\
\textbf{Veredicto:} \textbf{Correcto} (2/2 pts).

\subsection*{2. Sistemas de coordenadas (Longitud del hueso)}
\textbf{Procedimiento y resultado:} Identifica las coordenadas adecuadamente y escribe la fórmula de la distancia. Sustituye y resuelve los binomios al cuadrado sin errores ($\sqrt{36+64} = \sqrt{100}$), llegando a la distancia exacta de $10$. \\
\textbf{Veredicto:} \textbf{Correcto} (2/2 pts).

\subsection*{3. Vectores (Fuerza resultante)}
\textbf{Procedimiento y resultado:} Extrae de manera correcta las componentes rectangulares calculando por separado el eje $x$ ($45 + (-15) = 30$) y el eje $y$ ($25 + 35 = 60$). Sin embargo, a la hora de expresar la fuerza resultante final ($\vec{F}_R$) omite por completo los vectores unitarios ($\hat{i}$ y $\hat{j}$) y plantea una suma escalar plana ($30 + 60$), lo cual es un \textbf{error de notación y concepto} en el álgebra vectorial. \\
\textbf{Veredicto:} \textbf{Parcialmente correcto} (Cálculos numéricos correctos, error de notación en el resultado final) (1/2 pts).

\subsection*{4. Potencias de diez (Notación científica y prefijo)}
\textbf{Procedimiento y resultado:} Transforma la cifra decimal aplicando la regla de notación científica estricta ($1.25 \times 10^{-5}$) de forma totalmente correcta. No obstante, omite la segunda instrucción de la pregunta y no identifica ni nombra el prefijo del S.I. correspondiente. \\
\textbf{Veredicto:} \textbf{Incompleto / Parcialmente correcto} (1/2 pts).

\subsection*{5. Sistemas de unidades (Conversión de densidad)}
\textbf{Procedimiento y resultado:} La estudiante dejó esta pregunta completamente en blanco, sin plantear fórmulas ni intentos de cálculo en la segunda página. \\
\textbf{Veredicto:} \textbf{Incorrecto} (0/2 pts).

\vspace{0.5cm}
\section*{Comentarios Generales y Sugerencias}
La estudiante muestra un buen dominio aritmético inicial, resolviendo con mucha limpieza las primeras preguntas. 
\begin{itemize}
    \item Se le recomienda repasar con urgencia la notación de vectores, recordando que las componentes ortogonales ($\hat{i}, \hat{j}$) no se pueden sumar directamente como números simples (Pregunta 3).
    \item Es fundamental gestionar el tiempo y al menos intentar dejar un planteamiento en todos los ejercicios en lugar de dejarlos en blanco (Pregunta 5).
\end{itemize}

\end{document}